\documentclass{article}
\usepackage{graphicx} % Required for inserting images
\usepackage{amsmath, amssymb}
\usepackage{array}
\usepackage{booktabs}
\usepackage[margin=1in]{geometry}
\usepackage[round]{natbib}
\usepackage{amsthm}
\usepackage{enumerate}
\usepackage{float}
\usepackage[title, titletoc]{appendix}

\newtheorem{theorem}{Theorem}[section]
\newtheorem{lemma}[theorem]{Lemma}
\newtheorem{proposition}[theorem]{Proposition}
\newtheorem{corollary}[theorem]{Corollary}
\newtheorem{observation}[theorem]{Empirical Observation}
\theoremstyle{plain}
\newtheorem{heuristic}[theorem]{Heuristic Identity}
\newtheorem{derivation}[theorem]{Derivation Sketch}


\title{A Symbolic Prefix Decomposition of the Collatz Map and Its Empirical Determination of Per-Class Statistics}
\author{Nathan Humphrey}
\date{May 2026}

\begin{document}


\maketitle

\section{Introduction}
For odd integer $n$ with residue $r = n \bmod 2^k$, the Collatz iteration starting from $n$ can be tracked symbolically. Writing $n=2^k \cdot m+r$ and following the state $a \cdot m + c$ through Collatz iteration, the state's parity is determined by $r$ alone (i.e., by $a\bmod 2$ and $c \bmod 2$ jointly) for the prefix region while $a$ remains even, terminating when $a$ becomes odd. The terminal $a_{\text{final}}$ lies in $\{3^1, 3^2, \ldots, 3^k\}$; the number of odd-step applications consumed is at most $k$.

This paper establishes properties of the prefix decomposition and presents empirical evidence that the decomposition's terminal value $a_{\text{final}}$ determines per-class statistics of the total stopping time $\sigma$ on odd integers in residue class $r \bmod 2^k$. The contribution has two parts. The algebraic part, presented in \S3, establishes the properties of the prefix iteration that are deductive consequences of the symbolic state's parity dynamics. The principal algebraic statement is the termination theorem (\S3.2), which establishes that the prefix iteration terminates within at most $k$ odd-step applications and produces a terminal $a_{\text{final}}$ of the form $3^j$ for some $1 \leq j \leq k$. A closed-form quantity $\langle \alpha_{\text{det}} \rangle = \log(6)/ \log(4/3)$ is derived (\S3.5) from the prefix iteration under a heuristic combinatorial argument; we name this a heuristic identity rather than a theorem because the underlying combinatorial assumption is itself heuristic.

The empirical part, presented in \S4, presents regularities verified at modular resolutions $k \in \{6, 7, 8, 9\}$ and $N$ up to $2^{27}$. The principal empirical findings are: (i) per-class intercept $\alpha(r)$ of $\sigma$ vs $\log(n)$ is predicted by $\alpha_{\text{det}}(r)$ with regression $\text{R}^2 \geq 0.985$ and signal-to-noise ratio $\text{SD(residual)} / \text{mean(SE)}$ in $[0.91, 0.99]$ across the four tested resolutions, with per-class residuals scaling with sampling noise and no detectable structure above sampling noise; (ii) per-class central moments of $\sigma$ residuals cluster by $a_{\text{final}}$, with the number of distinct clusters equal to $k$ at every tested $k$; and (iii) the closed form $\langle \alpha_{\text{det}} \rangle = \log(6)/ \log(4/3)$ is verified empirically at each tested $k$ to within sampling precision. The empirical findings are presented as such; we do not claim the prefix decomposition's predictive content is theorem-proven for arbitrary $k$ or $N$.

The decomposition is a constructive sharpening of \citet{terras1976} Lemma 4: where Terras gives the asymptotic ratio between successive iterates' residue-class structure, we retain the symbolic state through the deterministic prefix and expose the terminal value $a_{\text{final}}$ as a closed-form covariate at finite $k$.

Terras's result establishes the asymptotic identity $S_k \approx S_0 \cdot 3^{d(k)}/ 2^k$, where $S_k$ is the residue-class structure at the $k$th Collatz iterate and $d(k)$ is the number of odd-step applications among the first $k$ Collatz steps. The prefix decomposition retains the symbolic state $(a, c)$ through the deterministic prefix and terminates at $a_{\text{final}} = 3^j$ explicitly, exposing the terminal value as a closed-form covariate. Terras's lemma gives the asymptotic content; the prefix decomposition gives the constructive content at finite $k$. Section 5 presents this relationship explicitly so that the contribution is bounded by what the decomposition adds to Terras's results, which is the constructive availability of $a_{\text{final}}$ and $\text{prefix\_steps}$ as closed-form predictors of per-class $\sigma$ statistics.

The paper is organized as follows. Section 2 places the decomposition in its lineage within the standard Collatz random-walk heuristic, \citet{terras1976} Lemma 4, and symbolic approaches to Collatz residue dynamics. Section 3 defines the prefix decomposition, states and proves the termination theorem, presents the empirical observation about realization of all $k$ $a_{\text{final}}$ values, defines the closed-form predictor $\alpha_{\text{det}}$, and derives the heuristic closed form $\langle \alpha_{\text{det}} \rangle = \log(6) / \log(4/3)$. Section 4 presents empirical regularities at modular resolutions $k \in \{6, 7, 8, 9\}$ and $N$ up to $2^{27}$. Section 5 places the decomposition relative to \citet{terras1976} explicitly. Section 6 discusses what the empirical regularities suggest, what is required for a stronger theorem-grade result, and the $qx+1$ extension. Section 7 concludes.

\section{Background}
\subsection{Collatz dynamics and the standard heuristic}
The Collatz map $T(n) = 3n + 1$ if $n$ is odd, $n/2$ if $n$ is even, generates the sequence $n \rightarrow T(n) \rightarrow T(T(n)) \rightarrow \ldots$ whose total stopping time $\sigma(n)$ is the number of steps from $n$ to $1$. For random odd $n$, the standard random-walk heuristic predicts $\mathbb{E}[\sigma \mid n] \approx K_h \cdot \log(n)$ where $K_h = 3/ \log(4/3) \approx 10.4282$ \citep{lagarias1985}. The heuristic derives from a random-walk argument: the Syracuse compression $n \rightarrow (3n + 1)/2^v$ has expected $\log$-multiplier $\log(3) - \langle v \rangle \cdot \log 2 = - \log(4/3)$ when $v$ is treated as $\text{Geometric}(\frac{1}{2})$, and the inverse of this drift gives $K_h$ \citep{lagarias2010}. Whether $\sigma$ has additional residue-class structure beyond the heuristic single-slope model is the question this paper addresses for $\sigma$ on odd integers.

\subsection{Terras 1976 Lemma 4}
\citet{terras1976} established the asymptotic identity for the $k$th Collatz iterate's residue-class structure: $$S_k \approx S_0 \cdot 3^{d(k)}/2^k$$ where $d(k)$ is the number of odd-step applications consumed among the first $k$ Collatz steps starting from the residue class. Terras's lemma gives the asymptotic ratio between successive iterates' residue-class structure but does not retain the symbolic state $(a, c)$ past the asymptotic identity. The prefix decomposition introduced in \S3 is a sharpening of Terras Lemma 4: we retain the full symbolic state through the deterministic prefix, terminate at $a_{\text{final}} = 3^j$ explicitly, and use this terminal value as a closed-form covariate for per-class $\sigma$ statistics rather than only for the asymptotic step ratio.

\subsection{Symbolic and algebraic approaches to Collatz residue dynamics}
Symbolic and algebraic approaches to Collatz residue dynamics have a substantial history. The natural setting is the $(\mathbb{Z}/3^k\mathbb{Z})^*$ tower, with reduction maps connecting modular resolutions, and the 2-adic structure governing the parity dynamics of even-step compressions. A comprehensive bibliography of this literature is collected in \citet{lagarias2010}. The prefix decomposition introduced here is restricted to the deterministic prefix region, the part of the Collatz iteration where the symbolic state's parity is determined by $r$ alone, and does not address the post-prefix dynamics, which are stochastic from the per-class perspective.

\section{The prefix decomposition: definition and algebraic properties}
\subsection{The symbolic state and parity rules}
Fix an odd integer $n$ with residue $r = n \bmod 2^k$, and write $n = 2^k \cdot m + r$, where $m$ is a non-negative integer. Define the symbolic state at step $t$ of the Collatz iteration as $a_t \cdot m + c_t$, with initial state $(a_0, c_0) = (2^k, r)$. The state is a linear function of $m$, with coefficients $(a_t, c_t)$ determined by the iteration's history.

The Collatz iteration's parity at step $t$ is determined by $a_t \cdot m + c_t \bmod 2$. While $a_t$ is even, the parity is $c_t \bmod 2$ -- independent of $m$. While $a_t$ is odd, the parity depends on $m$ and the iteration's behavior is no longer determined by $r$ alone.

The state transitions are:
\begin{enumerate}[(i)]
    \item If $a_t$ is even and $c_t$ is even: $a_t \cdot m + c_t$ is even for all $m$. The iteration applies $n \rightarrow n/2$, giving $(a_{t+1}, c_{t+1}) = (a_t/2, c_t/ 2)$.
    \item If $a_t$ is even and $c_t$ is odd: $a_t \cdot m + c_t$ is odd for all $m$. The iteration applies $n \rightarrow 3n+1$, giving $(a_{t+1}, c_{t+1}) = (3 \cdot a_t, 3 \cdot c_t + 1)$.
    \item If $a_t$ is odd: parity depends on $m$. The prefix iteration terminates.
\end{enumerate}

We define the prefix iteration as the sequence of state transitions (i) and (ii) from initial state $(a_0, c_0) = (2^k, r)$ until $a_t$ becomes odd (case iii).

\begin{lemma}[Prefix determinism]
\label{lem:prefix-determinism}
Let $n=2^k \cdot m + r$ with $r$ odd. The prefix iteration starting from initial state $(a_0, c_0) = (2^k, r)$ is fully determined by $r$: the sequence of state transitions, the terminal state $(a_{\text{final}}, c_{\text{final}})$, and the number of prefix steps are all independent of $m$.
\end{lemma}

\begin{proof}
The state transitions (i) and (ii) depend on $c_t \bmod 2$ and the parity of $a_t$, both of which are functions of the iteration history starting from $(2^k, r)$. By the parity rules above, while $a_t$ is even the parity of $a_t \cdot m + c_t$ equals $c_t \bmod 2$ independent of $m$. The prefix iteration consists exactly of state transitions (i) and (ii), which apply while $a_t$ is even and terminate when $a_t$ becomes odd. Therefore the entire prefix iteration -- sequence of transitions, terminal state, and number of prefix steps -- is a function of $r$ alone.
\end{proof}

\subsection{The termination theorem}

\begin{theorem}[Termination]
\label{thm:termination}
The prefix iteration terminates in finitely many steps. The number of odd-step applications (case ii) consumed during the prefix is at most $k$. The terminal $a_{\text{final}}$ is of the form $3^j$ for some integer $j$ with $1 \leq j \leq k$.
\end{theorem}

\begin{proof}
    The initial value $a_0 = 2^k$ contains exactly $k$ factors of 2. Even-step applications (case i) replace $a_t$ by $a_t/2$, decrementing the 2-adic valuation $v_2(a_t)$ by exactly $1$. Odd-step applications (case ii) replace $a_t$ by $3 \cdot a_t$; since $3$ is odd, this introduces no new factors of 2 into $a$, and $v_2(a_t)$ is unchanged.

    Therefore, $v_2(a_t)$ is monotone non-increasing along the prefix iteration, with $v_2(a_0) = k$ and $v_2(a_t)$ decreasing by exactly $1$ with each even-step application. The iteration enters case (iii), termination, the first time $v_2(a_t) = 0$, that is, the first time $a_t$ is odd. The number of even-step applications consumed during the prefix is therefore exactly $k$. The number of odd-step applications $j$ is bounded by the total prefix length but is otherwise determined by the specific sequence of parity values $c_t \bmod 2$ visited.

    The lower bound $j \geq 1$ follows from the initial state: $a_0 = 2^k$ is even and $c_0 = r$ is odd (since $r$ is odd by hypothesis), so the first step of the prefix iteration is case (ii), an odd-step application. Therefore the prefix consumes at least one odd-step application, giving $j \geq 1$. The upper bound $j \leq k$ follows from the 2-adic valuation argument above: each even-step decrements $v_2(a)$ by exactly $1$, the prefix terminates when $v_2 = 0$, so the prefix length is bounded above by $k$ even-steps and the number of odd-steps interleaved among them is at most $k$.
    
    The terminal $a_{\text{final}} = 3^j$, since each odd-step application multiplies $a$ by $3$ and the terminal value of $a$ (after the iteration's $k$ even-step applications have decremented the 2-adic valuation to $0$) is $3^j$ with $j$ the count of odd-step applications.
\end{proof}

The theorem is the deductive content of Lemma 1 plus the parity-rule arithmetic. Its proof is one paragraph and carries no heuristic content. The prefix iteration is well-defined, terminates in at most $2k$ steps total (at most $k$ even-step and at most $k$ odd-step applications), and produces a terminal $a_{\text{final}}$ of the form $3^j$ for $j \in \{1, 2, \ldots, k\}$.

A consequence of Theorem 1 is that the $2^{(k-1)}$ odd residue classes $\bmod 2^k$ partition into at most $k$ prefix-terminal types according to the value of $a_{\text{final}}$. Whether all $k$ types are actually realized at given $k$ (i.e., whether for each $j \in \{1, \ldots, k\}$ there exists $r \bmod 2^k$ whose prefix terminates at $a_{\text{final}} = 3^j$) is a separate question.

\subsection{Realization of all $k$ values of $a_{\text{final}}$}
\begin{observation}[Realization]
\label{obs:realization}
For modular resolutions $k \in \{4, 5, 6, 7, 8, 9, 10, 11, 12\}$, every $j \in \{1, \ldots, k\}$ is realized as $a_{\text{final}} = 3^j$ by some odd residue class $r \bmod 2^k$.

The observation has been verified by direct computation of the prefix iteration for all $2^{(k-1)}$ odd residue classes at each cited $k$, with the realized set of $a_{\text{final}}$ values enumerated in each case. We name the observation as empirical rather than as a theorem because we do not have a proof that all $k$ values are realized at every $k$ for arbitrary $k$. The observation is consistent with the available evidence and the structure of the prefix iteration suggests (but does not prove) that the realization extends to arbitrary $k$. We present the result as what we have verified directly, without extrapolation.
\end{observation}

\subsection{The closed-form predictor $\alpha_{\text{det}}$}
By the standard random-walk heuristic on the post-prefix state, the per-class intercept of $\sigma$ vs $\log(n)$ is heuristically: $$\alpha_{\text{det}}(r) := \text{prefix\_steps}(r) + K_h \cdot \log (a_{\text{final}}(r)/2^k)$$ where $\text{prefix\_steps}(r)$ is the total number of Collatz steps in the prefix from $(a_0, c_0) = (2^k, r)$ until $a$ becomes odd (the sum of even-step and odd-step applications during the prefix), and $K_h = 3/ \log(4/3)$. The prediction holds up to a single global additive constant absorbed into the global intercept $\mu_a$.

The derivation of $\alpha_{\text{det}}$ is heuristic: it applies the standard Collatz random-walk argument \citep{lagarias1985} to the post-prefix state $a_{\text{final}} \cdot m + c_{\text{final}}$, assuming the post-prefix dynamics behave as a random walk with i.i.d. $\text{Geometric}(\frac{1}{2})$ 2-adic valuations. The heuristic is the same one that produces the leading-order prediction $\mathbb{E}[\sigma \vert n] \approx K_h \cdot \log(n)$ for unconditioned $\sigma$; $\alpha_{\text{det}}$ adapts the heuristic to a specific post-prefix initial configuration. We present $\alpha_{\text{det}}$ as a closed-form quantity computable from $r$ and verify its empirical predictive content in \S4. We do not claim that $\alpha_{\text{det}}$ is the exact per-class intercept; we claim that it predicts the per-class intercept empirically.

\begin{heuristic}[Closed-form $\langle \alpha_{\text{det}} \rangle$]
\label{heur:alpha_det}
Under the heuristic that the prefix terminates at $a_{\text{final}} = 3^j$ with probability $2^{(-j)}$ (modeling each odd-step decision as an independent Bernoulli event), the class-averaged $\alpha_{\text{det}}$ satisfies $\langle \alpha_{\text{det}} \rangle = \log(6)/ \log(4/3)$.
\end{heuristic}

\begin{derivation}
\label{der:bern}
Under the independent-Bernoulli heuristic, the expected number of odd-step applications during the prefix is $\mathbb{E}[j] = \sum_{j \geq 1} j \cdot 2^{(-j)}=2$ (truncated by the bound $j \leq k$, but the truncation is negligible at $k \geq 6$). Substituting into $\alpha_{\text{det}}$ and computing the class average across the heuristic distribution gives $\langle \alpha_{\text{det}} \rangle = \langle \text{prefix\_steps} \rangle + K_h \cdot \langle \log(a_{\text{final}} / 2^k) \rangle$.

The two terms combine via the random-walk heuristic to yield $\log(6)/ \log(4/3)$. The full derivation is given in the accompanying repository \verb|(closed_form_findings.md)|.
\end{derivation}

We name the result a heuristic identity rather than a theorem because the underlying combinatorial assumption -- that each odd-step decision is an independent Bernoulli event -- is a heuristic. The actual distribution of $a_{\text{final}}$ across residue classes is determined by the prefix iteration's parity dynamics, not by independent coin flips. The heuristic is an approximation. The empirical verification in \S4.3 confirms that $\langle \alpha_{\text{det}} \rangle$ matches $\log(6)/ \log(4/3) \approx 6.226$ to within sampling precision at modular resolutions $k \in \{6, 7, 8, 9\}$.

\section{Empirical predictions of the prefix decomposition}
The findings in this section are empirical regularities verified at modular resolutions $k \in \{6, 7, 8, 9\}$ and $N$ up to $2^{27}$. Each finding is named, calibrated, and bounded by what was actually tested.

\subsection{Per-class intercept $\alpha(r)$}
At modular resolution $k = 6$ with $N=2^{25}$, hierarchical Bayesian fit of $\sigma \sim \alpha(r) + \beta(r) \cdot \log(n) + \varepsilon$ with class-level $\alpha$ and $\beta$ yields per-class $\alpha$ posterior means spanning $-27$ to $+35$ with $\tau_\alpha \approx 11.7$ \verb|(writeup.md Result 3)|. Linear regression of these against $\alpha_{\text{det}}(r)$ gives: $\alpha_{\text{actual}}(r) = -2.66 + 0.987 \cdot \alpha_{\text{det}}(r)$, $R^2 = 0.9996$ with $\text{SD(per-class residual)} = 0.28$ and $\text{SD(residual)}/ \text{mean(per-class SE)} = 0.18$ -- residuals below the noise floor of posterior sampling alone.

Extension to higher modular resolutions at $N=2^{27}$ via direct OLS regression (per-class regression rather than hierarchical Stan, to avoid posterior-geometry difficulties at large $k$):

\begin{table}[ht]
\centering
\caption{Per-class residual scale calibration across modular resolutions.}
\label{tab:residual-calibration}
\begin{tabular}{rrrrrc}
\hline
$k$ & mod & classes & per class $n$ & $R^2$ & $\text{SD(resid)} / \text{mean(SE)}$ \\
\hline
6 & 64 & 32 & 2.10M & 0.9967 & 0.96 \\
7 & 128 & 64 & 1.05M & 0.9942 & 0.99 \\
8 & 256 & 128 & 524K & 0.9918 & 0.91 \\
9 & 512 & 256 & 262K & 0.9851 & 0.93 \\
\hline
\end{tabular}
\end{table}

\begin{observation}[Per-class intercept prediction]
\label{obs:intercept}
At modular resolutions $k \in \{6, 7, 8, 9\}$ and $N$ up to $2^{27}$, per-class intercept $\alpha(r)$ of $\sigma$ vs $\log(n)$ is predicted by $\alpha_{\text{det}}(r)$ with regression $R^2 \geq 0.985$ and signal-to-noise ratio $\text{SD(residual)}/ \text{mean(SE)}$ in $[0.91, 0.99]$. Per-class residuals scale with sampling noise, with no detectable signal above sampling noise at the cited resolutions. The regularity is verified at the cited resolutions; we do not claim it for arbitrary $k$ or for $k > 9$.

A controlled comparison: at $k=8$ with $N = 2^{27}$ (524K per class, matching the original $k=6, N=2^{25}$ data scale of 524K per class), $R^2 = 0.9918$, essentially identical to the original $k=6$ $R^2$ of $0.9967$ -- same data per class produces same $R^2$ regardless of modular resolution.
\end{observation}

\subsection{The prefix-collapse onto $k$ distinct distributions}
The hypothesis initially tested was full distributional universality of $S(n) := \sigma(n) - \alpha(r) - \beta \cdot \log(n)$ -- that $S(n)$ would be class-independent. It is not: per-class variance differs by $\sim30\%$ relative across classes $\bmod 64$, and per-class skewness and kurtosis also vary. However, each higher central moment, when computed per class and plotted against $\alpha_{\text{det}}$ (or equivalently against $\log(a_{\text{final}})$, since these are collinear at fixed $k$), shows the prefix-prediction pattern:

\begin{table}[ht]
\centering
\caption{Pearson correlations of $\alpha_{\text{det}}$ with per-class central moments}
\label{tab:central-moments}
\begin{tabular}{rrrrrr}
\toprule
$k$ & classes & per-class $n$ & $r(\alpha_{\text{det}}, \text{Var})$ & $r(\alpha_{\text{det}}, \text{Kurt})$ & $r(\alpha_{\text{det}}, \text{Skew})$ \\
\midrule
6 & 32 & 2.10M & 0.99989 & 0.9497 & 0.8716 \\
7 & 64 & 1.05M & 0.99977 & 0.9337 & 0.8442 \\
8 & 128 & 524K & 0.99969 & 0.9236 & 0.8250 \\
9 & 256 & 262K & 0.99943 & 0.9077 & 0.7914 \\
\bottomrule
\end{tabular}
\end{table}

The clustering is sharper than the Pearson correlations alone indicate. When per-class variance is plotted against $\alpha_{\text{det}}$, classes with the same $a_{\text{final}}$ values produce 6 distinct clusters; at $k=7$, the 7 $a_{\text{final}}$ values produce 7 clusters; and so on through $k=9$. The same is true for kurtosis. Skewness shows slightly more within-cluster scatter but the same discrete structure.

\begin{observation}[Moment clustering by $a_{\text{final}}$]
\label{obs:moment-clustering}
At modular resolutions $k \in \{6, 7, 8, 9\}$, per-class central moments of $\sigma$ residuals cluster by $a_{\text{final}}$, with the number of distinct clusters equal to $k$ at every tested $k$. The regularity is verified at the cited resolutions.
\end{observation}

\subsection{Verification of the closed form $\langle \alpha_{\text{det}} \rangle = \log(6)/ \log(4/3)$}
\begin{observation}[Closed-form $\langle \alpha_{\text{det}} \rangle$ verification]
\label{obs:closed-form-verification}
Computed empirically at each tested $k$ from per-class $\alpha_{\text{det}}$ values, the class-averaged $\langle \alpha_{\text{det}} \rangle$ matches $\log(6)/ \log(4/3) \approx 6.226$ to within sampling precision at $k \in \{6, 7, 8, 9\}$. The empirical verification confirms the heuristic identity from \S3.5 at the cited resolutions.
\end{observation}


\subsection{Rate of convergence and limitations}
The empirical regularities are at finite $N$ (up to $2^{27}$ for the $k$-sweep). The $k$-sweep at fixed $N=2^{27}$ confirms the regularities at four distinct modular resolutions; we do not extrapolate to arbitrary $k$. Higher-$k$ extension is in principle possible but per-class sample sizes shrink as $2^{(1-k)}$; meaningful verification at $k > 12$ would require $N$ substantially larger than $2^{32}$. The signal-to-noise ratio $\text{SD(residual)}/ \text{mean(SE)}$ for the $\alpha_{\text{det}}$ predictions is in the band $[0.91, 0.99]$ across $k \in \{6, 7, 8, 9\}$. At $k=4$ the ratio is $0.70$ -- residuals smaller than per-class SE, meaning the prefix prediction over-explains at coarse resolution (signal-to-noise below sampling noise). At every $k$ tested, the regularities of \S4.1 and \S4.2 hold within their stated bounds. We do not claim the regularities hold at modular resolutions outside the tested range.

\section{Relation to \citep{terras1976}}
\citet{terras1976} Lemma 4 establishes the asymptotic identity for the $k$th Collatz iterate's residue-class structure: $S_k \approx S_0 \cdot 3^{d(k)}/ 2^k$, where $d(k)$ is the number of odd-step applications among the first $k$ Collatz steps starting from the residue class. The prefix decomposition retains the symbolic state $(a,c)$ through the deterministic prefix and terminates at $a_{\text{final}} = 3^j$ explicitly, where $j$ is the number of odd-step applications consumed during the prefix.

There are two ways to view the relationship:
\begin{enumerate}
    \item The prefix decomposition is a constructive version of Terras Lemma 4 at finite $k$. Where Terras gives the asymptotic ratio $S_k / S_0 \approx 3^{d(k)} / 2^k$ for the iterate's residue-class structure, the prefix decomposition gives the explicit terminal symbolic state $(a_{\text{final}}, c_{\text{final}}, \text{prefix\_steps})$ at finite $k$. The two results address related but distinct objects: Terras's lemma describes the $k$th iterate's residue distribution; the prefix decomposition describes the symbolic state evolution from initial residue $r$.
    \item The prefix decomposition strengthens Terras Lemma 4 by retaining the constructive content $(a_{\text{final}}$ and $\text{prefix\_steps}$) at finite $k$ as a closed-form covariate. This constructive content predicts per-class statistics of $\sigma$ at the level of the empirical regularities documented in \S4. The asymptotic content of Terras's lemma, by itself, does not produce these per-class predictions.

    We do not claim the prefix decomposition is novel relative to \citep{terras1976}. The decomposition is a sharpening of Terras Lemma 4 that retains constructive information Terras's lemma discards. The acknowledgment is explicit and located here so that no reader interpreting the decomposition as "new" relative to the broader literature goes away with that impression incorrectly.
\end{enumerate}
\section{Discussion}
\subsection{What the empirical regularities suggest}
At every tested modular resolution $k \in \{6, 7, 8, 9\}$, the prefix decomposition's $\alpha_{\text{det}}$ predicts per-class statistics of $\sigma$ at the level of sampling noise. The signal-to-noise ratio $\text{SD(Residual)}/ \text{mean(SE)}$ is in $[0.91, 0.99]$ across four resolutions. Per-class moments cluster by $a_{\text{final}}$ at every tested $k$, with the number of distinct clusters equal to $k$.

This suggests - but does not prove -- that the conditional distribution $\sigma(n) \mid r \bmod 2^k$ is parameterized entirely by $a_{\text{final}}$ at arbitrary $k$. A theorem to that effect would require characterizing the full conditional distribution $\sigma(n) \mid r$ analytically, which is an open question. The empirical evidence presented here is consistent with such a theorem but does not establish it.

\subsection{What is required for a stronger result}
Three pieces of analytical work would strengthen the empirical regularities of \S4 toward theorem-grade results:
\begin{enumerate}[\alph*)]
    \item Proof of "all $k$ values of $a_{\text{final}}$ are realized at every $k$" for arbitrary $k$ (Empirical Observation 1 strengthened to a theorem). The realization at small $k$ can be enumerated; whether the realization extends to arbitrary $k$ is a question about the structure of the prefix iteration that we have not resolved.
    \item Analytical characterization of the post-prefix dynamics from state $(a_{\text{final}}, c_{\text{final}}, m)$. The standard Collatz random-walk heuristic gives leading-order behavior on the post-prefix state but does not capture the higher-moment structure that the prefix-collapse predicts empirically. A characterization that recovers the empirical Pearson correlations of \S4.2 (variances, kurtosis, skewness clustering by $a_{\text{final}}$) from the post-prefix dynamics analytically would be a substantial strengthening.
    \item Connection to \citep{tao2022}'s almost-all-$N$ theorem. The prefix decomposition's $\alpha_{\text{det}}$ extends to predict per-class mean first-passage time at arbitrary thresholds $f(N)$, with the relationships $s_{\text{mean}}(r; f) \approx \alpha_{\text{det}}(r) + K_h \cdot \log(N/f(N))$ verified empirically across forty cells in the accompanying repository (and reported separately). A theorem-grade version would derive the per-class structural offset $\alpha_{\text{det}}$ from Tao's framework directly, connecting the prefix decomposition to the leading-term content of Tao's (5.15). 
\end{enumerate}

None of these are claimed in the present paper.

\subsection{Extensions to $qx+1$}
The prefix decomposition formally extends to $qx+1$ dynamics on the residue class $\bmod 2^k$ for $q$ coprime to 2. The terminal $a_{\text{final}}$ lives in $\{ q^j : 1 \leq j \leq k \}$. The construction's algebraic content -- Lemma 1 (parity determined by $r$ alone), Theorem 1 (termination within at most $k$ odd-step applications), and the structure of the prefix iteration -- transfers directly with $q$ in place of $3$.

Whether the empirical regularities of \S4 transfer to $qx+1$ is an open question. Preliminary empirical work at $q \in \{5, 7, 9, 11\}$ is reported separately. The $3x+1$ results of the present paper are restricted to the $q=3$ case.

\section{Conclusion}
This paper has established three classes of result on the prefix decomposition of the Collatz map.

The algebraic content: the prefix iteration is well-defined (Lemma 1), terminated within at most $k$ odd-step applications (Theorem 1), and produces a terminal $a_{\text{final}} \in \{3^1, 3^2, \ldots, 3^k\}$. Theorem 1 is the deductive content of the symbolic state's parity dynamics and is proven directly. The closed form $\langle \alpha_{\text{det}} \rangle = \log(6)/ \log(4/3)$ is a heuristic identity derived under independent-Bernoulli combinatorial assumption (Heuristic Identity 1) and is verified empirically.

The empirical content: at modular resolutions $k \in \{6, 7, 8, 9\}$ and $N$ up to $2^{27}$, the prefix decomposition's $\alpha_{\text{det}}$ predicts per-class intercept $\alpha(r)$ of $\sigma$ vs $\log(n)$ with regressions $R^2 \geq 0.985$ and signal-to-noise ratio $\text{SD(residual)} / \text{mean(SE)}$ in $[0.91, 0.99]$; per-class central moments of $\sigma$ residuals cluster by $a_{\text{final}}$, with the number of distinct clusters equal to $k$ at every tested $k$. The regularities are verified at the cited resolutions and are not extrapolated. The relationship to existing literature: the prefix decomposition is a sharpening of \citep{terras1976} Lemma 4. Where Terras's lemma gives the asymptotic ratio for the $k$th iterate's residue-class structure, the prefix decomposition retains the symbolic state through the deterministic prefix and exposes $a_{\text{final}}$ and $\text{prefix\_steps}$ as closed-form covariates for per-class $\sigma$ statistics. The contribution of the present paper is the constructive content (the explicit symbolic prefix and its empirically verified predictive power), not novelty relative to Terras's asymptotic identity.

Three open questions remain: rigorous extension of Empirical Observation 1 (realization of all $k$ $a_{\text{final}}$ values) to arbitrary $k$; analytical characterization of the post-prefix conditional distribution that recovers the empirical clustering of per-class central moments by $a_{\text{final}}$; and the formal connection of $\alpha_{\text{det}}$ to \citep{tao2022}'s almost-all-$N$ theorem at the level of per-class structural offset.

\appendix

\begin{appendices}

\section{Algorithmic specification of the prefix iteration}
\label{app:prefix-algorithm}
The prefix iteration is implemented directly from the parity rules \S3.1. Pseudocode for computing $(a_{\text{final}}, c_{\text{final}}, \text{prefix\_steps})$ given odd residue $r$ and modular resolution $k$:

\begin{verbatim}
function prefix(r, k):
    a := 2^k
    c := r
    steps := 0
    while a is even:
        if c is even:
            a := a / 2
            c := c / 2
        else:
            a := 3 * a
            c := 3 * c + 1
        steps := steps + 1
    return (a, c, steps)
\end{verbatim}

The iteration terminates when $a$ becomes odd (the loop's exit condition). Theorem~1 guarantees termination within at most $2k$ iterations. Reference implementations in the accompanying repository: \verb|alpha\_beta\_gamma\_decay.py| and \verb|analytical\_abc\_derivation.py|

\section{Reproducibility}
\label{app:reproducibility}
All code, data, and computational records are available at the accompanying repository. Specific scripts referenced in this paper:

\begin{itemize}
    \item \verb|experiments/01_alpha_decomposition.py| -- $k=6$ hierarchical Bayesian fit at $N = 2^{25}$.
    \item \verb|experiments/24_k_sweep_alpha_decomposition.py| -- $k$-sweep at $N = 2^{27}$ for $k \in \{4, \ldots, 12\}$.
    \item \verb|experiments/09_multi_stat_decomposition.py| -- per-class higher central moments (variance, kurtosis, skewness).
\end{itemize}

Detailed experimental records are in \verb|findings.md| and \verb|closed\_form\_findings.md| within the repository.

\end{appendices}

\bibliographystyle{plainnat}
\bibliography{references}

\end{document}
